\documentclass[11pt,notheorems,aspectratio=169,t]{beamer}

% Chinese package and font setup for Tao Slides Beamer.
% Microsoft YaHei mirrors the PowerPoint template. SimSun is a compile fallback.

\usepackage{fontspec}
\usefonttheme{professionalfonts}
\IfFileExists{/mnt/c/Windows/Fonts/msyh.ttc}{
  \setsansfont[
    Path=/mnt/c/Windows/Fonts/,
    BoldFont=msyhbd.ttc
  ]{msyh.ttc}
  \setmainfont[
    Path=/mnt/c/Windows/Fonts/,
    BoldFont=msyhbd.ttc
  ]{msyh.ttc}
}{
  \IfFontExistsTF{Microsoft YaHei}{
    \setsansfont{Microsoft YaHei}
    \setmainfont{Microsoft YaHei}
  }{
  \setsansfont{TeX Gyre Heros}
  \setmainfont{TeX Gyre Heros}
  }
}

\usepackage{xeCJK}
\IfFileExists{/mnt/c/Windows/Fonts/msyh.ttc}{
  \setCJKmainfont[
    Path=/mnt/c/Windows/Fonts/,
    BoldFont=msyhbd.ttc
  ]{msyh.ttc}
  \setCJKsansfont[
    Path=/mnt/c/Windows/Fonts/,
    BoldFont=msyhbd.ttc
  ]{msyh.ttc}
  \setCJKmonofont[
    Path=/mnt/c/Windows/Fonts/
  ]{msyh.ttc}
}{
  \IfFontExistsTF{Microsoft YaHei}{
    \setCJKmainfont[AutoFakeBold=2]{Microsoft YaHei}
    \setCJKsansfont[AutoFakeBold=2]{Microsoft YaHei}
    \setCJKmonofont{Microsoft YaHei}
  }{
  \setCJKmainfont[AutoFakeBold=2]{SimSun}
  \setCJKsansfont[AutoFakeBold=2]{SimSun}
  \setCJKmonofont{SimSun}
  }
}

\usepackage{styles/beamerthemetao}

\newcommand{\makepart}[1]{
  \part{#1}
  {
    \setbeamertemplate{footline}{}
    \begin{frame}
      \partpage
    \end{frame}
  }
}

\RequirePackage{booktabs}
\RequirePackage{colortbl}
\RequirePackage{ragged2e}
\RequirePackage{tabularx}
\RequirePackage{array}
\RequirePackage{multirow}
\RequirePackage{caption}
\RequirePackage{subcaption}
\RequirePackage{wrapfig}
\RequirePackage{pgfplots}
\RequirePackage{graphicx}
\RequirePackage{adjustbox}
\RequirePackage{environ}
\RequirePackage{textcomp}
\RequirePackage{amsmath}
\RequirePackage{amsthm}
\RequirePackage{mathtools}
\newcolumntype{Y}{>{\centering\arraybackslash}X}
\pgfplotsset{compat=1.18}

\DeclareMathOperator*{\argmax}{arg\,max}
\DeclareMathOperator*{\argmin}{arg\,min}
\setbeamertemplate{theorems}[numbered]

\theoremstyle{definition}
\newtheorem{fact}{事实}[section]
\newtheorem{examp}{示例}[section]

\theoremstyle{plain}
\newtheorem{definition}{定义}[section]
\newtheorem{proposition}{命题}
\newtheorem{theorem}{定理}
\newtheorem{assumption}{假设}

\providecommand{\H}{\mathscr{H}}
\providecommand{\E}{\mathbb{E}}


\title[标题]{标题}
\author{姓名}
\institute{单位}
\date{20XX年XX月XX日}

\begin{document}

{
\setbeamertemplate{footline}{}
\begin{frame}
  \titlepage
\end{frame}
}
\addtocounter{framenumber}{-1}

\begin{frame}{页标题}{副标题}
  \begin{itemize}
    \item 本模板用于创建简洁、统一的 Beamer 演示文稿。
    \item 页面风格包含白色背景、主题蓝标题、标题横线和右下角页码。
    \item 字体规范：页标题 36 pt，副标题 24 pt，正文 20 pt。
    \item 中文文本使用微软雅黑，数学公式使用 Computer Modern。
  \end{itemize}
\end{frame}

\begin{frame}{页标题}{副标题}
  \begin{itemize}
    \item 项目1
    \item 项目2
    \item 项目3
  \end{itemize}
\end{frame}

\begin{frame}{页标题}{副标题}
  \begin{itemize}
    \item 项目1
      \begin{itemize}
        \item 普通
        \item \emphasize{强调}
        \item \alert{重要}
      \end{itemize}
  \end{itemize}
\end{frame}

\begin{frame}{页标题}{副标题}
  \taobody{内容XX}
\end{frame}

\begin{frame}{页标题}{副标题}
  \taobody{%
    数学公式使用 Computer Modern 字体。

    \begin{align*}
      \frac{\partial n}{\partial t}
        &= \frac{1}{q}\nabla \cdot \mathbf{J}_n + G_n - R_n, \\
      \frac{\partial p}{\partial t}
        &= -\frac{1}{q}\nabla \cdot \mathbf{J}_p + G_p - R_p.
    \end{align*}
  }
\end{frame}

\end{document}
